\chapter{SCADA and HMI Standards}\label{ch:scada}
\section{Ignition Architecture}
Ignition is the preferred SCADA/HMI platform for new projects. The differences between a standard Ignition gateway and Ignition Edge Panel, and how the two are integrated, are covered in Appendix~\ref{app:edge-panel}. The project team should choose the Gateway, database, client, and redundancy arrangement based on system size, availability needs, performance, and recovery time. Operator workstations should normally run Perspective or Vision clients rather than host the production Gateway or database.

Record the tag providers, projects, databases, historian, alarm pipelines, identity provider, certificates, modules, backups, and important external dependencies. Use a development or test environment when practical instead of making untested changes directly in production.

\subsection[Version and Upgrades]{Version and Upgrades \newmark}\label{sec:ignition-version}
EFC runs Ignition~8.1 (Table~\ref{tab:software-versions}). New projects are built on the version the site's Gateway runs; a project is not built on a newer version than the Gateway that will host it. Modules that are not part of the standard Ignition install (third-party drivers, Sepasoft, and similar) need EFC approval because they add licence, support, and upgrade dependencies.
\future{An upgrade to Ignition~8.3 is planned. 8.3 keeps Gateway configuration as files on disk and adds a REST API, which changes how backups and source control work (Appendix~\ref{app:git}). Projects started before the upgrade should avoid features that do not carry over cleanly; the upgrade plan will list them.}

\section{Project and Tag Design}
Use shared projects, inheritance, UDTs, templates, and common scripts to avoid duplicating the same work across projects. Tag names and folders should follow the PLC equipment model and use consistent units, ranges, descriptions, quality handling, alarms, and history settings.

Keep scripts readable and reusable. Log errors that will help with troubleshooting, but avoid repetitive logging that hides useful information. Where a tag, expression, binding, or shared function can do the job cleanly, prefer it over copying scripts into many components.

\section[PLC Connections]{PLC Connections \newmark}\label{sec:ignition-plc-connections}
\begin{itemize}[itemsep=1pt]
  \item \textbf{One device connection per PLC}, named after the PLC (for example the project number and PLC name used in the file name, Section~\ref{sec:source-control}). Connection settings and the security mode are recorded in the system documentation.
  \item \textbf{Protocols} follow Section~\ref{sec:protocols}: OPC UA to M580 PLCs through the BMENUA0100, Modbus TCP to Productivity2000, and the Logix driver to the existing ControlLogix systems.
  \item \textbf{UDTs match PLC data types} member for member, so a PLC instance maps to one Ignition UDT instance with one parameter for its path (Section~\ref{sec:udt}).
  \item \textbf{Scan rates} (tag groups) are chosen for what the data is used for. Use a normal rate (typically 1~s) for most tags, a slower rate for maintenance counters and configuration, and a faster rate only where an operator or a calculation needs it. Leased or driven tag groups can reduce load for tags that are only shown on some screens.
  \item \textbf{Writes} from Ignition go only to command, setpoint, and configuration members (\texttt{Cmd\_}, \texttt{Prm\_}, \texttt{Cfg\_} and the current-type equivalents). Status and input members are read-only in the UDT. The PLC checks and clamps every written value (Section~\ref{sec:plc-setpoint-clamp}).
  \item \textbf{Bad quality} is shown on the screen, never hidden behind the last good value.
\end{itemize}

\section{HMI Philosophy}
Screens should make the current process condition easy to understand. Use color mainly to show abnormal conditions or items requiring attention. Do not rely on color alone; use text, symbols, and state indication as well. Navigation, faceplates, trends, units, decimal precision, manual controls, and alarm presentation should work the same way throughout the project.

Use a clear display hierarchy, normally including overview, area, unit, equipment-detail, and diagnostic screens. Operators should be able to see the active mode, command ownership, equipment state, alarms, interlocks, and bypasses before issuing a command. Use confirmation prompts for consequential actions without adding unnecessary prompts to routine work.

\section{Perspective Design}
Perspective screens follow the Perspective Screen Design Standard in Appendix~\ref{app:screen-design}. Perspective views should be tested at the screen sizes and aspect ratios expected in service. Use responsive containers where resizing is required. Docked views, popups, embedded views, parameters, and message handlers should have clear interfaces and should not depend on fragile component paths.

\paragraph{Script library.} Perspective screens call the project script library \texttt{efc.ui} (Ignition 8.1 and 8.3) for dynamic-window and command actions, so that every template behaves the same way. Template symbols call \texttt{openDynamic} on click and \texttt{dblCmd} on double-click in touch mode. Window close buttons call \texttt{closeDynamic}. The dynamic-window region calls \texttt{setDefaultDock} at startup to show its screen default view. \texttt{cmd} performs every command write after a role check. \texttt{bypass} and \texttt{unbypass} require Supervisor or higher and a reason, and write an audit entry. \texttt{ackAll}, \texttt{trend} and \texttt{level} provide alarm acknowledgement, trend display and the user's role level. Role levels and bypass rules follow Chapter~\ref{ch:scada} (\texttt{sec:ignition-roles}, \texttt{sec:hmi-bypass}).

Startup and navigation scripts should be safe to run more than once, finish quickly, and handle missing tags or bad-quality data without breaking the view. Reusable views should not assume that every optional parameter is present.

\section[Dynamic Windows]{Dynamic Windows \newmark}\label{sec:hmi-dynamic-windows}
Every equipment template opens a dynamic window: a panel on the right of the screen that shows the device's state, mode, commands, permissives, interlocks, alarms and history. Equipment of the same kind shares one dynamic window (Table~\ref{tab:hmi-dw-ids}), so a valve, motor or loop looks and works the same on every screen.

\begin{small}
\begin{longtable}{@{}P{3.6cm}P{12.6cm}@{}}
\caption{Dynamic Windows}\label{tab:hmi-dw-ids} \\
\toprule
\rowcolor{tblhead}\textbf{Window id} & \textbf{Used by} \\
\midrule
\endfirsthead
\multicolumn{2}{@{}l}{\small\textit{Table \thetable{} continued}} \\
\toprule
\rowcolor{tblhead}\textbf{Window id} & \textbf{Used by} \\
\midrule
\endhead
\texttt{dw\_ValveDiscrete} & On/off valves (\texttt{T\_DValve}) \\
\texttt{dw\_ValveAnalog} & Control valves (\texttt{T\_CValve}) \\
\texttt{dw\_Rotating} & Fixed-speed motors, pumps, fans, agitators and compressors (\texttt{T\_Motor}) \\
\texttt{dw\_VFD} & Variable-speed drives (\texttt{T\_VFD}) \\
\texttt{dw\_Heater} & Electric heaters (\texttt{T\_Heater}) \\
\texttt{dw\_HeatTransfer} & Heat exchangers (\texttt{T\_HeatEx}) \\
\texttt{dw\_ProcessValue} & Analog process values and indicators (\texttt{T\_AI}) \\
\texttt{dw\_Loop} & PID loops (\texttt{T\_PID}), including tuning \\
\texttt{dw\_Station} & Fill stations \\
\texttt{dw\_IOCard} & PLC I/O cards and racks (diagnostics) \\
\bottomrule
\end{longtable}
\end{small}

A template's \texttt{dockId} parameter overrides the default window chosen from its \texttt{tType} parameter, and \texttt{dwTagPath} overrides the bound tag. The \texttt{native} type opens the I/O card window; I/O card, rack and power-supply faceplates bind to the Modicon x80 device DDTs (decision 23).

I/O card, rack and power-supply faceplates bind directly to the Modicon x80 device DDTs from Control Expert, which are imported into Ignition as-is; there is no PCLib type for them (\texttt{T\_IOCard} and \texttt{T\_IORack} are removed), and the window id \texttt{dw\_IOCard} is kept.

\subsection[Size, Opening and Closing]{Size, Opening and Closing \newmark}\label{sec:hmi-dw-open}
\begin{itemize}[itemsep=1pt]
  \item \textbf{Size and position.} The dynamic window fills the full height between the header and the alarm bar, in the right-hand region of the shell: 400 $\times$ 860 at 1920 $\times$ 1080 (37.04\,vh wide, 79.63\,vh tall), and 800 $\times$ 1720 at 3840 $\times$ 2160 (Appendix~\ref{app:screen-design}).
  \item \textbf{Opening.} A single click or tap on the template symbol opens its dynamic window. The click target is the whole symbol bounding box, at least 44 design px in each direction.
  \item \textbf{Closing.} The close button returns the right-hand region to the screen's default view. Each screen view has a \texttt{defaultDock} view parameter that names it (normally the screen's trends), and \texttt{setDefaultDock} sets it on startup.
\end{itemize}

\subsection[Layout and Buttons]{Layout and Buttons \newmark}\label{sec:hmi-dw-layout}
Sections appear in this order; a section that does not apply to the device is left out:
\begin{enumerate}[itemsep=1pt]
  \item Header: tag, description and close button.
  \item State.
  \item Mode.
  \item Commands.
  \item Permissives.
  \item Interlocks.
  \item Alarms: active alarms are listed; inactive alarms are folded.
  \item History: an 8~h state strip and the most recent events.
  \item Maintenance and Configuration: collapsed by default and shown only to the roles in Table~\ref{tab:hmi-roles}. Loop tuning is done here, in \texttt{dw\_Loop}; there is no separate tuning popup.
\end{enumerate}
Buttons come in two sizes so that every dynamic window looks the same:
\begin{itemize}[itemsep=1pt]
  \item \textbf{Command tier} (mode AUTO /\allowbreak{} MAN /\allowbreak{} MAINT /\allowbreak{} OOS, OPEN /\allowbreak{} CLOSE, START /\allowbreak{} STOP, RESET, JOG, FWD /\allowbreak{} REV): one size, 72 $\times$ 30 design px, style \texttt{Button/Default}. The button for the current state is highlighted with \texttt{Button/Toggle/On}.
  \item \textbf{Secondary tier} (BYPASS, UNBYP, ACK, RESET LATCHED, RESET MAINT, SIMULATE, analog valve presets 0 /\allowbreak{} 25 /\allowbreak{} 50 /\allowbreak{} 75 /\allowbreak{} 100\,\%): one smaller size, style \texttt{Button/Small}.
  \item \texttt{Button/Danger} is used only for the confirm action in the bypass confirmation (Section~\ref{sec:hmi-bypass}).
\end{itemize}

\subsection[Touch Operation]{Touch Operation \newmark}\label{sec:hmi-touch}
In touch mode (template parameter \texttt{useButton = false}) a double tap on the symbol of a two-state device (\texttt{T\_DValve}, \texttt{T\_Motor}) sends the toggle command (open /\allowbreak{} close, start /\allowbreak{} stop). It is sent only when the device is in Manual and available, and is not interlocked, in local or faulted; otherwise the double tap does nothing. Each template has a \texttt{dblCmd} parameter (default on) to turn this off. Analog devices (control valves, drives, loops) take commands from the dynamic window only.

\subsection[Permissives, Interlocks and Bypass]{Permissives, Interlocks and Bypass \newmark}\label{sec:hmi-bypass}
The Permissives and Interlocks sections list each condition by name (\texttt{Cfg\_CondDesc} in \texttt{T\_Intlk}, Section~\ref{sec:pclib-t-intlk}) with its state, and mark the first-out, latched and bypassed conditions.
\begin{itemize}[itemsep=1pt]
  \item \textbf{Who.} Bypass and unbypass need the Supervisor role or higher. The BYPASS button is shown only for conditions configured as bypassable (\texttt{Cfg\_Bypassable}) and works only while the bypass authorisation \texttt{Inp\_BypassPerm} is on.
  \item \textbf{Confirmation.} BYPASS opens a confirmation (\texttt{Docks/Parts/BypassConfirm}). The user picks a duration from the presets 1~h, 2~h, 4~h and 8~h, limited to the group's maximum (\texttt{Cfg\_BypassMax}), and must enter a reason. The bypass is not applied without a reason.
  \item \textbf{Indication.} Each bypassed condition shows its own time left and an audit line: who bypassed it, their role, when, and the reason.
  \item \textbf{Logging.} Bypass and unbypass are written to the Ignition audit log (Section~\ref{sec:ignition-roles}).
  \item \textbf{Rule.} A bypass made this way from the HMI, on a condition configured as bypassable, is exempt from the second-person and approved MOC requirement in Section~\ref{sec:online-changes}. Safety-related signals are never configured as bypassable and cannot be bypassed in software. PLC forces and logic changes are not exempt.
\end{itemize}
\future{Named permissive and interlock groups are to be embedded in the device types: \texttt{Perm : T\_Intlk} and \texttt{Intlk : T\_Intlk} in \texttt{T\_DValve}, \texttt{T\_Motor} and \texttt{T\_VFD}, \texttt{Intlk : T\_Intlk} in \texttt{T\_CValve} and \texttt{T\_PID}, and \texttt{Prot : T\_Intlk} in equipment modules. The existing \texttt{Inp\_IntlkOK} and \texttt{Inp\_PermOK} are kept and ANDed with the group. \texttt{T\_Intlk} gains \texttt{Sts\_BypLeft : ARRAY[0..15] OF TIME} so the time left can be shown per condition; until then only the time left on the oldest bypass (\texttt{Sts\_BypassLeft}) is available.}

\subsection[Trends]{Trends \newmark}\label{sec:hmi-trends}
Trends in a dynamic window show the last 30~minutes by default. Every graph has an expand button. Expanding a graph keeps the dynamic window at 400 $\times$ 860 design px and the right-hand region never changes width: the graph fills the window below its header (about 400 $\times$ 760 design px) with range buttons 5~m, 30~m, 1~h, 8~h and 24~h, pause /\allowbreak{} live and a collapse button. The other sections of the window are hidden until the graph is collapsed.


\section{Alarms}
Configure Ignition alarms from the approved alarm list. Include the identifier, priority, message, cause, consequence, expected response, deadband, delay, shelving permissions, and any state-based suppression. Acknowledging an alarm does not clear the process condition that caused it. Alarm shelving needs the Supervisor role and alarm-limit changes the Engineer role, and both are audit-logged (Table~\ref{tab:hmi-roles}). \newmark

\textbf{Annunciation and acknowledgment.} Ignition annunciates alarms. For PLC data types that carry their own alarm state (the PCLib types, Appendix~\ref{app:data-types-future}), acknowledging the alarm in Ignition writes through to the alarm's \texttt{Cmd\_Ack} member in the PLC, so the PLC and SCADA states stay the same. Alarms are not acknowledged separately in two places. \newmark

\textbf{Priority.} Alarm priority uses Ignition's built-in priority levels. PLC data types that carry a priority use the same numbers, so a priority set in the PLC carries straight into Ignition without a lookup table: \newmark

\begin{small}
\begin{longtable}{@{}P{1.40cm}P{3.20cm}P{12.33cm}@{}}
\caption{Alarm Priority Numbering}\label{tab:alarm-priority} \\
\toprule
\rowcolor{tblhead}\textbf{Value} & \textbf{Ignition priority} & \textbf{Use} \\
\midrule
\endfirsthead
\toprule
\rowcolor{tblhead}\textbf{Value} & \textbf{Ignition priority} & \textbf{Use} \\
\midrule
\endhead
0 & Diagnostic & Equipment and system diagnostics for maintenance; not shown on the operator alarm summary by default. \\
1 & Low & Operator response needed, with time to respond and a minor consequence. \\
2 & Medium & Prompt response needed to avoid an upset or equipment damage. \\
3 & High & Immediate response needed to avoid a trip, release, or major equipment damage. \\
4 & Critical & Reserved for the most severe consequences; use sparingly. \\
\bottomrule
\end{longtable}
\end{small}
\note{PCLib v1.00 (Appendix~\ref{app:data-types-future}) still numbers \texttt{Cfg\_Priority} 1=Emergency to 4=Low. The generator will be changed to the numbering above at the next PCLib revision.}
Priority is assigned during alarm rationalization from the consequence and the time available to respond (Section~\ref{sec:alarm-mgmt} and Appendix~\ref{app:alarm-doc}), not chosen per tag by the programmer.

Test alarm journals and notification pipelines, including loss and recovery of the database or notification service. Review nuisance, chattering, standing, and frequently occurring alarms after startup.

\section{Historian and Data}
Historize data that has a defined operating, quality, maintenance, regulatory, or analysis purpose. Choose the sampling mode, deadband, timestamp source, and retention period to match that purpose. Document store-and-forward behavior and how bad or uncertain quality is handled.

Database connections should use managed credentials and only the permissions they need. Use encrypted connections when supported by the approved architecture.

\section{Security and Audit}
\subsection*{Individual accounts}\label{sec:ignition-accounts}
Every operator, engineer, and other user should have their own Ignition account and be logged in before using the system. Individual logins mean that actions such as setpoint changes, alarm acknowledgments, and bypasses are recorded against a person, and that each user gets only the access their role needs.
\note{Not currently implemented. This should be discussed in the near future, including how logins work on shared operator stations --- for example automatic logout after inactivity and a fast way to switch users at the station.}
\future{Ignition Perspective can use single sign-on (SSO) through an identity provider using OpenID Connect 1.0 or SAML 2.0, such as Microsoft Entra ID (Azure AD) or Okta. Users would log in with their existing company account rather than a separate Ignition password, and accounts would be added or removed centrally when people join or leave. This works for Perspective, Vision, the Designer, and the Gateway in both 8.1 and 8.3, and is worth looking into, possibly with the parent company's IT group.}

\subsection*{Roles and audit}\label{sec:ignition-roles}
Use the approved identity provider when practical and assign access by role. EFC uses six roles in a fixed order: Guest $<$ Operator $<$ Maintenance $<$ Supervisor $<$ Engineer $<$ Admin. Each role can do everything the roles below it can. The user's highest role is read once at login and held as a single session level, so a change to a user's roles takes effect at the next login. Screens compare that level with the level each action needs (Table~\ref{tab:hmi-roles}); controls that need a higher level than the user has are hidden. \newmark

\begin{small}
\begin{longtable}{@{}P{8.6cm}P{2.8cm}P{4.8cm}@{}}
\caption{Minimum Role for HMI Actions}\label{tab:hmi-roles} \\
\toprule
\rowcolor{tblhead}\textbf{Action} & \textbf{Minimum role} & \textbf{Audit log} \\
\midrule
\endfirsthead
\multicolumn{3}{@{}l}{\small\textit{Table \thetable{} continued}} \\
\toprule
\rowcolor{tblhead}\textbf{Action} & \textbf{Minimum role} & \textbf{Audit log} \\
\midrule
\endhead
View screens, trends and alarms & Guest & No \\
Mode AUTO /\allowbreak{} MAN, device commands in Manual, fault reset (Section~\ref{sec:pclib-modes}) & Operator & No \\
Alarm acknowledge, setpoint changes, reset of latched interlocks & \tbd & \tbd \\
Mode MAINT and OOS, JOG, simulation, maintenance resets (run time and counters) & Maintenance & No \\
Bypass and unbypass of permissives and interlocks (Section~\ref{sec:hmi-bypass}) & Supervisor & Yes, with reason \\
Alarm shelving & Supervisor & Yes \\
Alarm-limit changes & Engineer & Yes, old $\rightarrow$ new \\
Configuration and tuning (\texttt{Cfg\_} members, PID tuning) & Engineer & Yes, old $\rightarrow$ new \\
User and role administration & \tbd & Yes \\
\bottomrule
\end{longtable}
\end{small}

Bypass, unbypass, shelving, alarm-limit, configuration and tuning changes are written to the Ignition audit log with the old and new value, the user and role, the time, and the reason where one is entered. Also audit setpoint changes, recipe changes, and security administration.
\note{Open: minimum role for alarm acknowledge, setpoint changes and reset of latched interlocks, and which actions are reserved for Admin (TBD above).}

OPC UA connections should use trusted certificates, a supported security policy (Basic256Sha256, sign and encrypt), named accounts with the minimum permissions needed, and no anonymous access (Section~\ref{sec:protocols}). Remove default accounts and unused connections before turnover.

\section{Backup, Deployment, and Recovery}
Create a Gateway backup before and after significant production changes and after commissioning. Record the Ignition version, installed modules, licenses, database dependencies, certificates, identity-provider settings, DNS/NTP needs, and restore steps.

For critical systems, demonstrate that the backup can be restored and that the project can be rolled back if a deployment fails. Keep the accepted backup with the project turnover package.
